% !TEX root = ../main.tex
\chapter{関連研究}\label{chap:related}
% ===== 下書き(Claude、11/22)。レオが読んで理解 → 確認の問い → 直す =====
% 注意:各研究の中身は、題目・要旨・先行発表の予稿の記述に基づく。表の「要確認」は原文を読んで確かめる
% 11/22 確認済み:Moore ら(USENIX の本文)、Tang ら(要旨、レオが貼付)、Zanini ら(本文 PDF、レオ)、坂本ら(J-GLOBAL の抄録、レオ)、笹倉ら(本文 PDF、レオ)

本章では、本研究に関わる研究を、データセンタの熱の再循環、データセンタの冷却の予測制御、プロセッサの予測に基づく熱管理、不確かさを考慮したマージン、エッジ環境の熱管理の順に述べる。最後に、本研究の位置づけを示す。

\section{データセンタの熱の再循環と熱を考慮した配置}\label{sec:rel-recirc}
データセンタでは、サーバの排熱が別のサーバの吸気側に回り込む熱の再循環が、冷却の効率を下げる原因として知られている。Moore ら\cite{moore2005cool}は、熱の再循環を考慮して負荷をどのサーバに割り当てるかを決める方法を提案し、気流のシミュレーション(CFD)で、負荷の置き場所を考えない場合と比べて冷却のコストを最大で半分にできることを示した。Tang ら\cite{tang2008sensor}は、熱の再循環を線形のモデルで表し、吸気温度の最大値が最も小さくなるようにタスクを割り当てる問題として定式化した。小規模なデータセンタのシミュレーションで、他の方法より吸気温度を 2〜5\,℃ 下げ、冷却のエネルギーを 20〜30\,\% 減らせることを示した。

これらの研究は、本研究の空間熱干渉と同じく、ある計算機の排熱が別の計算機の温度を上げる現象を扱っている。ただし、再循環の大きさは事前に求めた係数として与えられ、操作するのは負荷の配置である。本研究では、隣のノードからの熱の影響を時間とともに変わる状態として推定し、操作するのはノードごとのファンである。また、本研究の環境(開いたラック)では、この影響はほぼ生じなかった(第\ref{chap:preliminary}章)。

\section{データセンタの冷却の予測制御}\label{sec:rel-dc-mpc}
データセンタの冷却設備の制御にも、予測が用いられている。Ogura ら\cite{ogura2018mpc}は、コールドアイルを囲ったデータセンタを対象に、空調の給気温度と稼働させるサーバの台数をモデル予測制御で決め、全体の消費電力を小さくする方法を提案し、シミュレーションで評価した。モデル予測制御は、モデルで将来の状態を予測し、制約を守りながら目的を最もよくする操作を、時刻ごとに求め直す方法である。

また、笹倉ら\cite{sasakura2020}は、実際のサーバルームで収集した空調機と各ラックの消費電力、吸気温度のデータから、状態空間モデルと機械学習、熱収支の式を組み合わせて、各ラックの吸気温度を 30 分先まで予測した。状態空間モデルでは予測の 95\,\% 信頼区間の上下限も求め、実測が予測から外れる幅を見積もれることを示した。データセンタでは、ラックの吸気温度を一定以下に保つことが利用者との契約(SLA)で定められることが多く、温度の予測とその不確かさが運用に必要とされている。ただし、この研究は予測までであり、予測を用いて冷却を操作することは扱っていない。

部屋の空調のように、操作してから効果が現れるまでに時間がかかる冷却では、温度を見てから操作したのでは間に合わない場合がある。本研究は、ノードごとのファンという効きの速い冷却を対象とし、予測を用いる方式としきい値方式を実機で比べた点が異なる。また、冷却の効きの遅さを変えたシミュレーションで、どの程度遅い冷却から予測が有効になるかを調べた(第\ref{chap:eval}章)。
% 要文献:空調の応答の時間(数十秒〜分)を実測で示した文献

\section{プロセッサの予測に基づく熱管理}\label{sec:rel-dtm}
プロセッサの温度を一定の範囲に保つ技術は、動的熱管理と呼ばれる。Skadron ら\cite{skadron2003temperature}は、プロセッサの各部の温度を熱 RC モデルで表し、温度が上がりすぎないように動作を制限する方法を検討した。Zanini ら\cite{zanini2009multicore}は、マルチコアプロセッサの熱 RC モデルを用いたモデル予測制御で、各コアの動作周波数を決め、温度の上限を守りながら性能を保つ方法を提案した。

Zanini らのモデルは、チップを小さな区画に分け、隣り合う区画の間の熱の伝わりを含む熱 RC の回路で表しており、あるコアの熱が隣のコアに伝わる影響も含んでいる。ただし、その係数はチップの構造から事前に計算したものであり、評価は Niagara-1 の構成を模したシミュレーションで行っている。

これらの研究では、温度を下げる手段として、動作周波数を下げるなどプロセッサの処理を抑える方法を用いる。この方法は温度を下げる代わりに性能を下げる。本研究では、性能を保ったままファンで冷やし、ファンの稼働を減らすことを目指す。また、これらの研究はモデルの係数が正確に分かっていることを前提としており、隣からの熱の影響を観測から推定することや、予測の不確かさを扱うことはしていない。

\section{不確かさを考慮したマージン}\label{sec:rel-margin}
予測には誤差があるため、予測に基づいて上限を守るには、誤差を見込んだ余裕(マージン)が必要である。確率的モデル予測制御\cite{mesbah2016smcp}では、制約が破られる確率を一定以下に抑えるという形で不確かさを扱い、予測の分布に応じて制約に余裕を持たせる。計算機の資源管理では、Gmach ら\cite{gmach2007capacity}が負荷の予測に基づく容量の計画を提案している。坂本ら\cite{sakamoto2018kbse}は、物理マシンの資源を超えて仮想マシンを割り当てる(オーバーコミット)環境で、仮想マシンの CPU 使用率の予測に持たせるマージンを、固定ではなく各仮想マシンの変動の大きさに合わせて動的に変える方法を提案した。変動の大きい仮想マシンには大きなマージンを、安定した仮想マシンには小さなマージンを与えることで、仮想マシンの移動(ライブマイグレーション)に用いたときに SLA の違反が減り、資源の利用効率が上がることを実験で示した。

坂本らの方法は、過去の予測の変動の大きさからマージンを決める。これに対し本研究では、カルマンフィルタの誤差共分散から求めた予測の標準偏差 $\sigma$ でマージンを決める。誤差共分散は、隣のノードの温度が急に変わったときなど、モデルの上で不確かさが増える状況を反映できる。ただし、本研究の環境では干渉がほぼないため、$\sigma$ は一定となり、この利点は現れなかった(第\ref{chap:eval}章)。

\section{エッジ環境の熱管理}\label{sec:rel-edge}
エッジの機器は、空調のない場所に置かれることが多く、周囲の温度の影響を受けやすい。Ahmadi ら\cite{edgeengine2023}は、エッジでの推論について、周囲の温度によって性能が変わることを踏まえ、温度を考慮して推論の設定を最適化する枠組みを提案した。Raspberry Pi を束ねたクラスタは、安価な計算基盤として教育や研究に用いられている\cite{iridispi2014}。
% 要確認:EdgeEngine の操作の対象(推論の設定・DVFS など)を原文で確かめる

\section{本研究の位置づけ}\label{sec:rel-position}
表\ref{tab:rel-compare}に、関連研究と本研究の違いをまとめる。本研究の特徴は次の3つである。

\begin{itemize}
  \item 操作量が、性能を下げない冷却(ノードごとのファン)であり、温度の上限を守ることに加えて、ファンの稼働を減らすことを目的とする
  \item 予測の不確かさを、カルマンフィルタの誤差共分散から求め、マージンに用いる
  \item 実機のファンを動かし、閉ループで評価する。そのうえで、予測が有効になる条件(熱干渉の有無、冷却の効きの速さ)を示す
\end{itemize}

\begin{table}[tb]
  \centering
  \caption{関連研究との比較}
  \label{tab:rel-compare}
  \small
  \begin{tabular}{p{0.15\linewidth}p{0.12\linewidth}p{0.15\linewidth}p{0.13\linewidth}p{0.13\linewidth}p{0.12\linewidth}}
    \toprule
    研究 & 対象 & 操作するもの & 予測の方法 & マージン & 評価 \\
    \midrule
    Moore ら\cite{moore2005cool} & データセンタ & 負荷の配置 & 再循環を考えた配置の方針 & なし & シミュレーション(CFD) \\
    Tang ら\cite{tang2008sensor} & データセンタ & タスクの割り当て & 線形の再循環モデル & なし & シミュレーション \\
    笹倉ら\cite{sasakura2020} & サーバルーム & なし(予測のみ) & 状態空間モデル + 機械学習 & 95\,\% 信頼区間 & 実測データ \\
    Ogura ら\cite{ogura2018mpc} & データセンタ & 空調の給気温度、サーバの台数 & モデル予測制御 & なし & シミュレーション \\
    Zanini ら\cite{zanini2009multicore} & マルチコア & 動作周波数・電圧 & 熱 RC + モデル予測制御 & なし & シミュレーション \\
    坂本ら\cite{sakamoto2018kbse} & 仮想マシン & 仮想マシンの配置 & CPU 使用率の予測 & 変動の大きさに応じて動的 & 実験 \\
    本研究 & マイクロクラスタ & ファン(ON/OFF) & 熱 RC + カルマンフィルタ & 誤差共分散 & 実機(閉ループ) \\
    \bottomrule
  \end{tabular}
\end{table}
